% ═══════════════════════════════════════════════════════════════════════════
%  فصل ۱۰ — مداخلات معاصر: ۵۰ سال اخیر در سراسر جهان
% ═══════════════════════════════════════════════════════════════════════════
\chapter{مداخلات معاصر: ۵۰ سال اخیر در سراسر جهان}
\label{ch:contemporary}

\begin{tcolorbox}[mybox, title={چکیده‌ی فصل}]
این فصل به بررسی تفصیلی مهم‌ترین مداخلات خارجی و نمونه‌های نجات‌طلبی در نیم‌قرن اخیر (۱۹۷۵–۲۰۲۴) می‌پردازد. از جنگ خلیج فارس (۱۹۹۱) و سومالی (۱۹۹۲) تا اوکراین (۲۰۲۲)، هر مورد با مدل شش‌وجهی تحلیل و در طیف نتایج جای‌گذاری می‌شود. هدف آن است که از مقایسه‌ی تطبیقی این موارد، الگوهای تکرارشونده و درس‌های قابل‌تعمیم استخراج گردد.
\end{tcolorbox}

% ─────────────────────────────────────────────────────────────────────────
\section{مقدمه: دوران پس از جنگ سرد و «نظم نوین جهانی»}
\label{sec:ch10-intro}
% ─────────────────────────────────────────────────────────────────────────

پایان جنگ سرد (۱۹۸۹–۱۹۹۱) تحول بنیادینی در الگوی مداخلات خارجی ایجاد کرد. در دوران جنگ سرد، مداخلات عمدتاً در چارچوب رقابت ابرقدرتی (امریکا-شوروی) صورت می‌گرفت و هر مداخله با «وتوی» طرف مقابل در شورای امنیت مواجه بود. پس از فروپاشی شوروی، \thinker{جورج بوش پدر} از «نظم نوین جهانی» (\lr{New World Order}) سخن گفت — نظمی که در آن شورای امنیت بتواند آزادانه‌تر عمل کند.%
\footnote{\lr{Freedman, Lawrence. \textit{A Choice of Enemies: America Confronts the Middle East}. PublicAffairs, 2008, pp.~210--235.}}

\thinker{فرانسیس فوکویاما} (\lr{Francis Fukuyama}) در «پایان تاریخ» (۱۹۹۲) ادعا کرد که دموکراسی لیبرال «آخرین مرحله‌ی تکامل سیاسی بشر» است — و بنابراین مداخله برای «دموکراسی‌سازی» مشروعیت جهان‌شمول دارد.%
\footnote{\lr{Fukuyama, Francis. \textit{The End of History and the Last Man}. Free Press, 1992.}}
اما تجربه‌ی سه دهه‌ی اخیر نشان داد که واقعیت بسیار پیچیده‌تر است.

\begin{figure}[htbp]
\centering
\begin{tikzpicture}[
  x=0.017\linewidth,
  every node/.style={font=\scriptsize},
]

% خط زمان
\draw[line width=2.5pt, PrimaryMid] (0,0) -- (130,0);

% رویدادها — بالا
\node[circle, fill=EraContemp, minimum size=7pt, inner sep=0pt] at (0,0) {};
\node[above=0.6cm, text width=2cm, align=center] at (0,0)
  {\rl{خلیج فارس}\\1991};

\node[circle, fill=AccentRed, minimum size=7pt, inner sep=0pt] at (15,0) {};
\node[above=0.6cm, text width=2cm, align=center] at (15,0)
  {\rl{سومالی}\\1992};

\node[circle, fill=AccentRed, minimum size=7pt, inner sep=0pt] at (25,0) {};
\node[above=0.6cm, text width=2cm, align=center] at (25,0)
  {\rl{رواندا}\\1994};

\node[circle, fill=AccentAmber, minimum size=7pt, inner sep=0pt] at (40,0) {};
\node[above=0.6cm, text width=2cm, align=center] at (40,0)
  {\rl{کوزوو}\\1999};

\node[circle, fill=EraContemp, minimum size=7pt, inner sep=0pt] at (55,0) {};
\node[above=0.6cm, text width=2cm, align=center] at (55,0)
  {\rl{افغانستان}\\2001};

% رویدادها — پایین
\node[circle, fill=AccentRed, minimum size=7pt, inner sep=0pt] at (70,0) {};
\node[below=0.6cm, text width=2cm, align=center] at (70,0)
  {\rl{عراق}\\2003};

\node[circle, fill=AccentAmber, minimum size=7pt, inner sep=0pt] at (85,0) {};
\node[below=0.6cm, text width=2cm, align=center] at (85,0)
  {\rl{لیبی}\\2011};

\node[circle, fill=AccentRed, minimum size=7pt, inner sep=0pt] at (95,0) {};
\node[below=0.6cm, text width=2cm, align=center] at (95,0)
  {\rl{سوریه}\\2011--};

\node[circle, fill=AccentRed, minimum size=7pt, inner sep=0pt] at (110,0) {};
\node[below=0.6cm, text width=2cm, align=center] at (110,0)
  {\rl{یمن}\\2015--};

\node[circle, fill=EraPostmod, minimum size=7pt, inner sep=0pt] at (125,0) {};
\node[below=0.6cm, text width=2cm, align=center] at (125,0)
  {\rl{اوکراین}\\2022--};

\end{tikzpicture}
\caption{خط زمانی مداخلات معاصر (۱۹۹۱–۲۰۲۴)}
\label{fig:ch10-timeline}
\end{figure}

% ─────────────────────────────────────────────────────────────────────────
\section{جنگ خلیج فارس (۱۹۹۱): مداخله‌ی «مشروع»}
\label{sec:ch10-gulf-war}
% ─────────────────────────────────────────────────────────────────────────

\subsection{زمینه و ماهیت}

حمله‌ی \enemy{صدام حسین} به کویت (اوت ۱۹۹۰) و اخراج ارتش عراق توسط ائتلاف بین‌المللی تحت رهبری امریکا (ژانویه–فوریه ۱۹۹۱) اغلب به‌عنوان \concept{موفق‌ترین مداخله‌ی نظامی پس از جنگ سرد} معرفی می‌شود. دلایل:

\begin{itemize}[nosep, rightmargin=1em]
  \item مجوز صریح شورای امنیت (قطعنامه ۶۷۸).
  \item ائتلاف ۳۴ کشوری (چندجانبه‌گرایی واقعی).
  \item هدف محدود و مشخص: آزادی کویت (نه تغییر رژیم عراق).
  \item استراتژی خروج: پس از آزادی کویت، ائتلاف متوقف شد.%
  \footnote{\lr{Freedman, Lawrence, and Efraim Karsh. \textit{The Gulf Conflict, 1990--1991}. Faber, 1993.}}
\end{itemize}

\subsection{اما: شورش شیعیان و کردها}

نکته‌ی تلخ جنگ خلیج فارس، سرنوشت \concept{شورش شعبانیه} بود. در مارس ۱۹۹۱، شیعیان جنوب و کردهای شمال عراق — با تشویق ضمنی بوش پدر — علیه صدام شوریدند. اما ائتلاف مداخله نکرد و صدام شورش را به‌شدت سرکوب کرد (حدود ۱۰۰٬۰۰۰ کشته).%
\footnote{\lr{Tripp, Charles. \textit{A History of Iraq}. 3rd ed. Cambridge UP, 2007, pp.~246--260.}}

\begin{tcolorbox}[enemybox, title={درس تاریخی: وعده‌ی مداخله بدون عمل}]
بوش پدر در فوریه ۱۹۹۱ گفت: «مردم عراق باید امور را به دست خود بگیرند و دیکتاتور صدام حسین را کنار بگذارند.»%
\footnote{\lr{Bush, George H.W. Public Statement, 15 February 1991.}}

اما وقتی مردم عراق دقیقاً همین کار را کردند، امریکا مداخله نکرد. دلایل:
\begin{itemize}[nosep, rightmargin=1em]
  \item ترس از تجزیه‌ی عراق و تقویت ایران.
  \item فقدان مجوز شورای امنیت برای تغییر رژیم.
  \item تمایل عربستان و مصر به حفظ تمامیت عراق.
\end{itemize}

\textbf{نتیجه:} بار دیگر الگوی «بی‌وفایی حامی خارجی» تکرار شد. شیعیان و کردهای عراق از این تجربه درس گرفتند: هنگام حمله‌ی ۲۰۰۳، کردها با احتیاط بیشتری عمل کردند و شیعیان ابتدا بی‌اعتماد بودند.
\end{tcolorbox}

% ─────────────────────────────────────────────────────────────────────────
\section{سومالی (۱۹۹۲–۱۹۹۵): مداخله‌ی انسان‌دوستانه‌ی شکست‌خورده}
\label{sec:ch10-somalia}
% ─────────────────────────────────────────────────────────────────────────

\subsection{زمینه و تحلیل}

فروپاشی دولت سومالی (۱۹۹۱) و قحطی گسترده، سازمان ملل و سپس امریکا را به مداخله‌ی نظامی واداشت. عملیات \lr{UNOSOM I} (۱۹۹۲) و سپس \lr{UNITAF} (عملیات «بازگرداندن امید») با هدف تأمین امنیت کمک‌رسانی آغاز شد.%
\footnote{\lr{Clarke, Walter, and Jeffrey Herbst, eds. \textit{Learning from Somalia: The Lessons of Armed Humanitarian Intervention}. Westview, 1997.}}

\begin{tcolorbox}[casebox, title={تحلیل شش‌وجهی: سومالی (۱۹۹۲)}]
\begin{itemize}[nosep, rightmargin=1em]
  \item \textbf{ساختاری:} دولت شکست‌خورده؛ هیچ «حکومتی» وجود نداشت که مداخله را «دعوت» کند. جنگ‌سالاران حاکم بودند.
  \item \textbf{حقوقی:} قطعنامه ۷۹۴ شورای امنیت (فصل هفتم). نخستین بار شورای امنیت «بحران انسان‌دوستانه» را تهدیدی علیه صلح و امنیت بین‌المللی خواند.%
  \footnote{\lr{UN Security Council Resolution 794, 3 December 1992.}}
  \item \textbf{اقتصادی:} سومالی فاقد منابع استراتژیک بود — بنابراین انگیزه‌ی اقتصادی ناچیز.
  \item \textbf{نظامی:} نبرد موگادیشو (اکتبر ۱۹۹۳) — کشته‌شدن ۱۸ سرباز امریکایی — موجب عقب‌نشینی شد.%
  \footnote{\lr{Bowden, Mark. \textit{Black Hawk Down}. Atlantic Monthly, 1999.}}
  \item \textbf{روان‌شناختی:} «اثر سی‌ان‌ان» (\lr{CNN Effect}) — تصاویر قحطی افکار عمومی غرب را بسیج کرد، اما تصاویر سرباز کشته‌شده‌ی امریکایی آن را فرو ریخت.%
  \footnote{\lr{Robinson, Piers. \textit{The CNN Effect: The Myth of News, Foreign Policy and Intervention}. Routledge, 2002.}}
  \item \textbf{ژئوپولیتیک:} سومالی اهمیت ژئوپولیتیک کمی داشت. عقب‌نشینی هزینه‌ی سیاسی کمتری داشت.
  \item \textbf{نتیجه (طیف):} \textbf{E}. شکست کامل؛ سومالی تا ۲۰۲۴ هنوز فاقد دولت مرکزی مؤثر است.
\end{itemize}
\end{tcolorbox}

% ─────────────────────────────────────────────────────────────────────────
\section{رواندا (۱۹۹۴): عدم مداخله و نسل‌کشی}
\label{sec:ch10-rwanda}
% ─────────────────────────────────────────────────────────────────────────

\subsection{فاجعه‌ی عدم اقدام}

\enemy{نسل‌کشی رواندا} (آوریل–ژوئیه ۱۹۹۴) — کشتار حدود ۸۰۰٬۰۰۰ توتسی و هوتوی میانه‌رو در ۱۰۰ روز — مهم‌ترین نمونه‌ی \concept{شکست عدم مداخله} در تاریخ معاصر است.%
\footnote{\lr{Dallaire, Roméo. \textit{Shake Hands with the Devil: The Failure of Humanity in Rwanda}. Random House, 2003.}}

\begin{histQuote}{ژنرال رومئو دالر، فرمانده نیروی سازمان ملل در رواندا}
«من با ۵٬۰۰۰ سرباز مجهز می‌توانستم نسل‌کشی را متوقف کنم. شورای امنیت نه‌تنها نیرو نفرستاد بلکه نیروهای موجود را نیز کاهش داد ... ما رواندایی‌ها را رها کردیم.»%
\footnote{\lr{Dallaire, 2003, p.~375.}}
\end{histQuote}

\subsection{چرا مداخله نشد؟}

\begin{tcolorbox}[failbox, title={عوامل عدم مداخله در رواندا}]
\begin{enumerate}[nosep, label=\textbf{\arabic*.}]
  \item \textbf{اثر سومالی:} شکست سومالی (شش ماه قبل) امریکا را از هرگونه مداخله در آفریقا بازداشت. دستورالعمل \lr{PDD-25} (مه ۱۹۹۴) شرایط سختی برای مداخله تعیین کرد — درست هنگام نسل‌کشی.%
  \footnote{\lr{Power, Samantha. \textit{A Problem from Hell: America and the Age of Genocide}. Basic Books, 2002, pp.~329--389.}}
  \item \textbf{فقدان منافع ژئوپولیتیک:} رواندا نه نفت داشت، نه اهمیت استراتژیک.
  \item \textbf{فقدان اطلاعات؟} نه — سازمان ملل و سیا از برنامه‌ریزی نسل‌کشی مطلع بودند.%
  \footnote{\lr{Melvern, Linda. \textit{Conspiracy to Murder: The Rwandan Genocide}. Verso, 2004.}}
  \item \textbf{اجتناب از واژه «نسل‌کشی»:} امریکا عمداً از به‌کاربردن واژه «نسل‌کشی» خودداری کرد تا از تعهد حقوقی مداخله فرار کند (کنوانسیون ۱۹۴۸).%
  \footnote{\lr{Power, 2002, pp.~358--364.}}
  \item \textbf{نقش فرانسه:} فرانسه — متحد حکومت هوتو — عملاً مانع مداخله شد و حتی نیروهایش (\lr{Opération Turquoise}) متهم به حمایت از عاملان نسل‌کشی هستند.%
  \footnote{\lr{Prunier, Gérard. \textit{The Rwanda Crisis: History of a Genocide}. 2nd ed. Columbia UP, 1997.}}
\end{enumerate}
\end{tcolorbox}

\begin{tcolorbox}[defbox, title={پارادوکس رواندا}]
رواندا نشان داد که «عدم مداخله» نیز می‌تواند فاجعه‌بار باشد. این فاجعه محرک اصلی تدوین دکترین \concept{مسئولیت حمایت} (\lr{R2P}) شد — همان‌گونه که در فصل ۲ بحث شد. اما آیا وجود \lr{R2P} می‌تواند رواندای بعدی را جلوگیری کند؟ تجربه‌ی سوریه (بخش ۱۰.۸) نشان می‌دهد که پاسخ منفی است.
\end{tcolorbox}

% ─────────────────────────────────────────────────────────────────────────
\section{کوزوو (۱۹۹۹): مداخله بدون مجوز}
\label{sec:ch10-kosovo}
% ─────────────────────────────────────────────────────────────────────────

\subsection{زمینه و مشروعیت}

بمباران ۷۸روزه‌ی ناتو علیه صربستان (مارس–ژوئن ۱۹۹۹) برای توقف \concept{پاک‌سازی قومی} آلبانیایی‌تبارهای کوزوو توسط \enemy{میلوشویچ}، مهم‌ترین نمونه‌ی «مداخله‌ی انسان‌دوستانه بدون مجوز شورای امنیت» است.%
\footnote{\lr{Independent International Commission on Kosovo. \textit{The Kosovo Report}. Oxford UP, 2000.}}

\begin{legalNote}{معضل حقوقی کوزوو}
کمیسیون مستقل بین‌المللی کوزوو نتیجه‌گرفت که مداخله‌ی ناتو \textbf{«غیرقانونی اما مشروع»} (\lr{illegal but legitimate}) بود:
\begin{itemize}[nosep, rightmargin=1em]
  \item \textbf{غیرقانونی:} فاقد مجوز شورای امنیت (وتوی روسیه و چین).
  \item \textbf{مشروع:} پاک‌سازی قومی گسترده و فوریت بشردوستانه.
\end{itemize}
این فرمول‌بندی یکی از بحث‌برانگیزترین مفاهیم حقوق بین‌الملل معاصر شد.%
\footnote{\lr{Kosovo Report, 2000, p.~4.}}
\end{legalNote}

\begin{tcolorbox}[casebox, title={تحلیل شش‌وجهی: کوزوو (۱۹۹۹)}]
\begin{itemize}[nosep, rightmargin=1em]
  \item \textbf{ساختاری:} سرکوب سیستماتیک آلبانیایی‌ها توسط صربستان؛ ارتش آزادی‌بخش کوزوو (\lr{KLA}) فعال اما ضعیف.
  \item \textbf{حقوقی:} بدون مجوز شورای امنیت. «غیرقانونی اما مشروع.»
  \item \textbf{اقتصادی:} کوزوو فقیر و فاقد منابع استراتژیک. انگیزه‌ی اقتصادی ناچیز.
  \item \textbf{نظامی:} بمباران هوایی بدون نیروی زمینی. تلفات غیرنظامی قابل‌توجه (حدود ۵۰۰ کشته‌ی غیرنظامی صرب).%
  \footnote{\lr{Human Rights Watch. "Civilian Deaths in the NATO Air Campaign." February 2000.}}
  \item \textbf{روان‌شناختی:} حافظه‌ی سربرنیتسا (۱۹۹۵) و رواندا (۱۹۹۴) فشار اخلاقی قوی ایجاد کرده بود.
  \item \textbf{ژئوپولیتیک:} ناتو نیاز داشت پس از جنگ سرد «هدف جدید» بیابد. کوزوو آن هدف شد. روسیه مخالف بود اما ضعیف‌تر از آن بود که جلوگیری کند.%
  \footnote{\lr{Judah, Tim. \textit{Kosovo: War and Revenge}. Yale UP, 2000.}}
  \item \textbf{نتیجه (طیف):} \textbf{B}. استقلال کوزوو (۲۰۰۸) با حمایت غرب؛ اما شناسایی ناقص بین‌المللی و وابستگی شدید به ناتو و اتحادیه‌ی اروپا.
\end{itemize}
\end{tcolorbox}

% ─────────────────────────────────────────────────────────────────────────
\section{افغانستان (۲۰۰۱–۲۰۲۱): بیست سال مداخله و فروپاشی}
\label{sec:ch10-afghanistan}
% ─────────────────────────────────────────────────────────────────────────

\subsection{از «جنگ با تروریسم» تا «دولت‌سازی»}

حمله‌ی امریکا به افغانستان (اکتبر ۲۰۰۱) پس از حملات ۱۱ سپتامبر، با هدف نابودی القاعده و سرنگونی طالبان آغاز شد. هدف اولیه نظامی بود نه «نجات مردم افغانستان». اما به‌تدریج، هدف به \concept{دولت‌سازی} (\lr{State-Building}) و \concept{دموکراسی‌سازی} تغییر یافت.%
\footnote{\lr{Rashid, Ahmed. \textit{Descent into Chaos}. Viking, 2008.}}

\subsection{سقوط کابل (آگوست ۲۰۲۱): شکست نهایی}

بیست سال مداخله‌ی نظامی، بیش از ۲ تریلیون دلار هزینه، و حدود ۱۷۰٬۰۰۰ کشته (از جمله ۴۷٬۰۰۰ غیرنظامی افغان) با سقوط فوری دولت افغانستان و بازگشت طالبان در ۱۵ آگوست ۲۰۲۱ پایان یافت.%
\footnote{\lr{Costs of War Project. Watson Institute, Brown University. "Afghan War." Updated 2021.}}

\begin{tcolorbox}[casebox, title={تحلیل شش‌وجهی: افغانستان (۲۰۰۱–۲۰۲۱)}]
\begin{itemize}[nosep, rightmargin=1em]
  \item \textbf{ساختاری:} دولت مرکزی ضعیف و فاسد (غنی، کرزی)؛ وفاداری‌ها قومی-قبیله‌ای؛ نهادهای مدرن سطحی و بدون ریشه.
  \item \textbf{حقوقی:} مجوز شورای امنیت (قطعنامه ۱۳۶۸ و ۱۳۷۳) و حق دفاع مشروع (ماده ۵۱ منشور).
  \item \textbf{اقتصادی:} افغانستان فاقد منابع معدنی قابل‌دسترس (لیتیوم بهره‌برداری‌نشده)؛ اقتصاد مبتنی بر کمک خارجی و مواد مخدر.%
  \footnote{\lr{Special Inspector General for Afghanistan Reconstruction (SIGAR). \textit{What We Need to Learn}. August 2021.}}
  \item \textbf{نظامی:} برتری نظامی مطلق امریکا اما ناتوانی در مهار شورش (\lr{counterinsurgency}). طالبان از پناهگاه‌های پاکستان بهره برد.
  \item \textbf{روان‌شناختی:} «آرزواندیشی» غربی: باور به امکان تبدیل افغانستان به «دموکراسی» در عرض یک دهه — بدون توجه به واقعیت‌های قبیله‌ای و فرهنگی. از سوی دیگر، بخشی از شهرنشینان افغان (زنان، تحصیل‌کرده‌ها) واقعاً از مداخله «نجات‌طلبانه» استقبال کردند.%
  \footnote{\lr{Maley, William. \textit{The Afghanistan Wars}. 3rd ed. Palgrave, 2021.}}
  \item \textbf{ژئوپولیتیک:} بازی بزرگ جدید: پاکستان (حمایت از طالبان)، ایران (نفوذ در غرب)، روسیه و چین (مخالفت ضمنی).
  \item \textbf{نتیجه (طیف):} \textbf{D–E}. بازگشت طالبان؛ فروپاشی فوری دولت؛ از دست‌رفتن دستاوردهای بیست‌ساله (به‌ویژه برای زنان)؛ بی‌اعتمادی جهانی به «دولت‌سازی» غربی.
\end{itemize}
\end{tcolorbox}

% ─────────────────────────────────────────────────────────────────────────
\section{عراق (۲۰۰۳): مداخله‌ی پیش‌گیرانه و فاجعه‌ی دولت‌سازی}
\label{sec:ch10-iraq}
% ─────────────────────────────────────────────────────────────────────────

\subsection{از «سلاح‌های کشتار جمعی» تا «دموکراسی‌سازی»}

حمله‌ی امریکا و بریتانیا به عراق (مارس ۲۰۰۳) بدون مجوز شورای امنیت و با ادعای نادرست وجود «سلاح‌های کشتار جمعی» (\lr{WMD}) صورت گرفت. \thinker{هانس بلیکس} (\lr{Hans Blix})، رئیس بازرسان سازمان ملل، بعدها تأیید کرد که ادعای \lr{WMD} بی‌اساس بود.%
\footnote{\lr{Blix, Hans. \textit{Disarming Iraq}. Pantheon, 2004.}}

\subsection{نجات‌طلبی بخشی از عراقی‌ها}

بر خلاف افغانستان، در مورد عراق \textbf{بخشی از جامعه‌ی عراقی واقعاً از مداخله استقبال کرد:}
\begin{itemize}[nosep, rightmargin=1em]
  \item \textbf{کردها:} پس از تجربه‌ی انفال و منطقه‌ی پرواز ممنوع (۱۹۹۱–۲۰۰۳)، کردها قوی‌ترین حامیان مداخله بودند.
  \item \textbf{شیعیان:} پس از تجربه‌ی تلخ ۱۹۹۱ (شورش سرکوب‌شده)، بخشی از شیعیان — به‌ویژه احزاب تبعیدی (مجلس اعلا، حزب الدعوه) — از مداخله حمایت کردند. اما توده‌ی شیعیان بی‌اعتماد بودند.%
  \footnote{\lr{Allawi, Ali A. \textit{The Occupation of Iraq: Winning the War, Losing the Peace}. Yale UP, 2007.}}
  \item \textbf{احمد چلبی و «کنگره‌ی ملی عراق»:} چلبی فعالانه از حمله‌ی امریکا حمایت و حتی اطلاعات جعلی درباره‌ی \lr{WMD} ارائه کرد.%
  \footnote{\lr{Ricks, Thomas E. \textit{Fiasco: The American Military Adventure in Iraq}. Penguin, 2006.}}
\end{itemize}

\begin{tcolorbox}[enemybox, title={نتایج فاجعه‌بار عراق}]
\begin{itemize}[nosep, rightmargin=1em]
  \item \textbf{کشته‌ها:} بیش از ۲۰۰٬۰۰۰ غیرنظامی عراقی (تخمین‌های مختلف).%
  \footnote{\lr{Iraq Body Count. \url{https://www.iraqbodycount.org/}. Accessed 2024.}}
  \item \textbf{آوارگان:} بیش از ۴ میلیون.
  \item \textbf{هزینه:} بیش از ۲ تریلیون دلار برای امریکا.
  \item \textbf{ظهور داعش:} انحلال ارتش عراق و بعث‌زدایی، زمینه‌ساز ظهور داعش (۲۰۱۴) شد.%
  \footnote{\lr{Cockburn, Patrick. \textit{The Rise of Islamic State}. Verso, 2015.}}
  \item \textbf{تقویت ایران:} مداخله‌ی امریکا عملاً ایران را به قدرتمندترین بازیگر خارجی در عراق تبدیل کرد.
  \item \textbf{نتیجه (طیف):} \textbf{C–D}. تغییر رژیم اما بی‌ثباتی مزمن، فساد، خشونت فرقه‌ای، وابستگی ساختاری.
\end{itemize}
\end{tcolorbox}

% ─────────────────────────────────────────────────────────────────────────
\section{لیبی (۲۰۱۱): اِعمال \texorpdfstring{\lr{R2P}}{R2P} و هرج‌ومرج پس از آن}
\label{sec:ch10-libya}
% ─────────────────────────────────────────────────────────────────────────

\subsection{مداخله بر مبنای مسئولیت حمایت}

مداخله‌ی ناتو در لیبی (مارس–اکتبر ۲۰۱۱) بر اساس قطعنامه‌ی ۱۹۷۳ شورای امنیت و تحت عنوان «حمایت از غیرنظامیان» انجام شد — نخستین اِعمال عملی دکترین \lr{R2P} با مجوز شورای امنیت.%
\footnote{\lr{Bellamy, Alex J., and Paul D. Williams. "The New Politics of Protection? Côte d'Ivoire, Libya and the Responsibility to Protect." \textit{International Affairs} 87, no.~4 (2011): 825--850.}}

اما ناتو سریعاً از «حمایت از غیرنظامیان» به «تغییر رژیم» حرکت کرد — سقوط و قتل \enemy{معمر قذافی} (اکتبر ۲۰۱۱).

\subsection{هرج‌ومرج پس از قذافی}

\begin{tcolorbox}[failbox, title={لیبی پس از ۲۰۱۱: دولت شکست‌خورده}]
\begin{itemize}[nosep, rightmargin=1em]
  \item دو (و گاه سه) دولت رقیب.
  \item جنگ داخلی مستمر.
  \item گسترش سلاح به ساحل و شمال آفریقا (بحران مالی ۲۰۱۲).
  \item مرکز قاچاق انسان و مهاجرت غیرقانونی به اروپا.%
  \footnote{\lr{Kuperman, Alan J. "Obama's Libya Debacle." \textit{Foreign Affairs} 94, no.~2 (2015): 66--77.}}
  \item \textbf{اوباما} بعدها لیبی را «بدترین اشتباه» دوران ریاست‌جمهوری‌اش خواند.%
  \footnote{\lr{Goldberg, Jeffrey. "The Obama Doctrine." \textit{The Atlantic}, April 2016.}}
\end{itemize}
\end{tcolorbox}

\begin{tcolorbox}[defbox, title={لیبی و بحران \lr{R2P}}]
لیبی \textbf{آخرین مورد اِعمال \lr{R2P} با مجوز شورای امنیت} بود. دلیل: روسیه و چین — که در رأی‌گیری قطعنامه ۱۹۷۳ «ممتنع» مانده بودند — پس از مشاهده‌ی تغییر رژیم، اعلام کردند فریب خورده‌اند و قول دادند هرگز مجوز مشابهی ندهند. نتیجه: \textbf{بن‌بست شورای امنیت در سوریه.}%
\footnote{\lr{Zifcak, 2012.}}
\end{tcolorbox}

% ─────────────────────────────────────────────────────────────────────────
\section{سوریه (۲۰۱۱–کنون): شکست هنجاری}
\label{sec:ch10-syria}
% ─────────────────────────────────────────────────────────────────────────

\subsection{از اعتراضات مسالمت‌آمیز تا جنگ داخلی تمام‌عیار}

بحران سوریه با اعتراضات مسالمت‌آمیز (مارس ۲۰۱۱) علیه \enemy{بشار اسد} آغاز شد و به‌سرعت به جنگ داخلی و سپس «جنگ نیابتی» تبدیل گردید — با مداخله‌ی مستقیم روسیه، ایران، ترکیه، امریکا و ده‌ها گروه مسلح.%
\footnote{\lr{Lesch, David W. \textit{Syria: The Fall of the House of Assad}. Yale UP, 2012.}}

\subsection{سوریه از منظر نجات‌طلبی}

بحران سوریه \textbf{پیچیده‌ترین مورد نجات‌طلبی معاصر} است زیرا تقریباً \textbf{همه‌ی بازیگران هم‌زمان «نجات‌طلب» و «مداخله‌گر»} بودند:

\begin{longtable}{
  >{\raggedleft\arraybackslash\bfseries}p{2.5cm}
  >{\raggedleft\arraybackslash}p{3cm}
  >{\raggedleft\arraybackslash}p{3cm}
  >{\raggedleft\arraybackslash}p{4cm}}
\caption{بازیگران بحران سوریه و رابطه‌ی نجات‌طلبی/مداخله}\label{tab:ch10-syria-actors}\\
\toprule
\rowcolor{PrimaryDeep}
\textcolor{white}{\textbf{بازیگر}} &
\textcolor{white}{\textbf{از چه کسی «نجات» می‌خواهد؟}} &
\textcolor{white}{\textbf{از چه کسی کمک می‌خواهد؟}} &
\textcolor{white}{\textbf{نتیجه}} \\
\midrule
\endfirsthead
\toprule
\rowcolor{PrimaryDeep}
\textcolor{white}{\textbf{بازیگر}} &
\textcolor{white}{\textbf{از چه کسی «نجات» می‌خواهد؟}} &
\textcolor{white}{\textbf{از چه کسی کمک می‌خواهد؟}} &
\textcolor{white}{\textbf{نتیجه}} \\
\midrule
\endhead
\bottomrule
\endlastfoot

اپوزیسیون سنّی &
رژیم اسد &
ترکیه، عربستان، قطر، امریکا &
تکه‌تکه شدن، شکست نظامی%
\footnote{\lr{Lister, Charles R. \textit{The Syrian Jihad}. Oxford UP, 2015.}} \\
\rowcolor{NeutralLight}
کردها (روژاوا) &
اسد + داعش + ترکیه &
امریکا (ائتلاف ضدداعش) &
خودمختاری موقت؛ رهاشدن توسط امریکا (۲۰۱۹) \\
دولت اسد &
اپوزیسیون + داعش &
روسیه، ایران، حزب‌الله &
بقا (اما بر ویرانه‌ها)%
\footnote{\lr{Hokayem, Emile. \textit{Syria's Uprising and the Fracturing of the Levant}. Routledge, 2013.}} \\
\rowcolor{NeutralLight}
ایزدی‌ها &
داعش (نسل‌کشی) &
کردها + ائتلاف (بمباران هوایی) &
بقا اما آوارگی گسترده%
\footnote{\lr{UN Human Rights Council. "They Came to Destroy." A/HRC/32/CRP.2, 2016.}} \\

\end{longtable}

\begin{tcolorbox}[enemybox, title={آمار فاجعه‌ی سوریه (تا ۲۰۲۳)}]
\begin{itemize}[nosep, rightmargin=1em]
  \item بیش از ۵۰۰٬۰۰۰ کشته.
  \item بیش از ۶.۸ میلیون آواره‌ی خارجی.
  \item بیش از ۶.۹ میلیون آواره‌ی داخلی.
  \item تخریب بیش از ۵۰٪ زیرساخت‌های کشور.%
  \footnote{\lr{Syrian Observatory for Human Rights, 2023. UN OCHA Syria Crisis Overview, 2023.}}
  \item \textbf{نتیجه (طیف):} \textbf{E–F}. فاجعه‌ی تمام‌عیار. هیچ‌یک از بازیگران به هدف خود نرسید مگر روسیه (حفظ اسد) و ایران (حفظ نفوذ).
\end{itemize}
\end{tcolorbox}

% ─────────────────────────────────────────────────────────────────────────
\section{اوکراین (۲۰۱۴–کنون): نجات‌طلبی، الحاق و جنگ تمام‌عیار}
\label{sec:ch10-ukraine}
% ─────────────────────────────────────────────────────────────────────────

\subsection{کریمه (۲۰۱۴): نجات‌طلبی یا الحاق؟}

الحاق کریمه توسط روسیه (مارس ۲۰۱۴) از سوی مسکو با ادعای «حمایت از روس‌زبانان» و «رفراندوم» مشروعیت‌بخشی شد. اما رفراندوم تحت اشغال نظامی و بدون نظارت بین‌المللی برگزار شد و مجمع عمومی آن را محکوم کرد.%
\footnote{\lr{Sakwa, Richard. \textit{Frontline Ukraine: Crisis in the Borderlands}. I.B. Tauris, 2015.}}

\subsection{تهاجم ۲۰۲۲ و نجات‌طلبی اوکراینی}

تهاجم تمام‌عیار روسیه به اوکراین (فوریه ۲۰۲۲) پیچیده‌ترین بحران امنیتی اروپا از ۱۹۴۵ به بعد است. از منظر نجات‌طلبی:

\begin{tcolorbox}[casebox, title={اوکراین: نجات‌طلبی دوسویه}]
\begin{itemize}[nosep, rightmargin=1em]
  \item \textbf{نجات‌طلبی روسیه:} پوتین ادعا کرد که «روس‌زبانان دونباس» از «نسل‌کشی» توسط اوکراین رنج می‌برند و روسیه وظیفه‌ی «نجات» آنها را دارد — ادعایی که سازمان ملل و اکثر نهادهای بین‌المللی آن را رد کردند.%
  \footnote{\lr{UN General Assembly Resolution ES-11/1, 2 March 2022.}}

  \item \textbf{نجات‌طلبی اوکراین:} اوکراین از ناتو، اتحادیه‌ی اروپا و جهان درخواست کمک نظامی، مالی و انسان‌دوستانه کرد — و دریافت کرد. این \textbf{یکی از معدود نمونه‌های نجات‌طلبی مؤثر} معاصر است.

  \item \textbf{تفاوت کلیدی با موارد قبلی:} اوکراین دارای \concept{دولت مشروع}، \concept{ارتش قوی}، \concept{هویت ملی منسجم} و \concept{حمایت گسترده‌ی بین‌المللی} بود — شرایطی که در سومالی، لیبی و افغانستان وجود نداشت.

  \item \textbf{اما:} وابستگی ساختاری به کمک غربی، خطر «خستگی حامیان» (\lr{donor fatigue})، و عدم‌قطعیت درباره‌ی پایان جنگ — همه نشان‌دهنده‌ی خطرات بلندمدت نجات‌طلبی حتی در «بهترین» شرایط است.
\end{itemize}
\end{tcolorbox}

% ─────────────────────────────────────────────────────────────────────────
\section{جدول تطبیقی جامع: مداخلات معاصر}
\label{sec:ch10-comprehensive-table}
% ─────────────────────────────────────────────────────────────────────────

\begin{landscape}
\begin{table}[htbp]
\centering
\caption{جدول تطبیقی جامع مداخلات معاصر (۱۹۹۱–۲۰۲۴)}\label{tab:ch10-comprehensive}
\begin{adjustbox}{max width=\linewidth}
\begin{tabular}{
  >{\raggedleft\arraybackslash\bfseries\small}p{1.8cm}
  >{\raggedleft\arraybackslash\small}p{1cm}
  >{\raggedleft\arraybackslash\small}p{1.5cm}
  >{\raggedleft\arraybackslash\small}p{1.3cm}
  >{\raggedleft\arraybackslash\small}p{1.2cm}
  >{\raggedleft\arraybackslash\small}p{1.2cm}
  >{\raggedleft\arraybackslash\small}p{1.5cm}
  >{\raggedleft\arraybackslash\small}p{1.5cm}
  >{\raggedleft\arraybackslash\small}p{1.5cm}
  >{\raggedleft\arraybackslash\small}p{1cm}}
\toprule
\rowcolor{PrimaryDeep}
\textcolor{white}{مورد} &
\textcolor{white}{سال} &
\textcolor{white}{مداخله‌گر} &
\textcolor{white}{مجوز \lr{SC}} &
\textcolor{white}{دعوت؟} &
\textcolor{white}{نوع} &
\textcolor{white}{هزینه انسانی} &
\textcolor{white}{نتیجه فوری} &
\textcolor{white}{نتیجه بلندمدت} &
\textcolor{white}{طیف} \\
\midrule

خلیج فارس &
۱۹۹۱ &
ائتلاف ۳۴ کشور &
بله &
کویت &
نظامی &
متوسط &
آزادی کویت &
ثبات کویت &
\cellcolor{BgSuccess}A \\
\rowcolor{NeutralLight}
سومالی &
۱۹۹۲ &
امریکا/سازمان ملل &
بله &
— &
انسان‌دوستانه &
محدود &
شکست &
دولت شکست‌خورده &
\cellcolor{BgFail}E \\
رواندا &
۱۹۹۴ &
(عدم مداخله) &
— &
— &
— &
۸۰۰٬۰۰۰ &
نسل‌کشی &
بازسازی &
\cellcolor{BgFail}F \\
\rowcolor{NeutralLight}
کوزوو &
۱۹۹۹ &
ناتو &
خیر &
\lr{KLA} &
بمباران &
محدود &
توقف پاک‌سازی &
استقلال شکننده &
\cellcolor{BgSuccess!60}B \\
افغانستان &
۲۰۰۱ &
امریکا/ناتو &
بله &
ائتلاف شمال &
نظامی &
۱۷۰K &
سقوط طالبان &
بازگشت طالبان &
\cellcolor{BgFail!80}D–E \\
\rowcolor{NeutralLight}
عراق &
۲۰۰۳ &
امریکا/بریتانیا &
خیر &
اپوزیسیون &
نظامی &
۲۰۰K+ &
سقوط صدام &
بی‌ثباتی، داعش &
\cellcolor{AccentAmber!30}C–D \\
لیبی &
۲۰۱۱ &
ناتو &
بله &
\lr{NTC} &
بمباران &
متوسط &
سقوط قذافی &
هرج‌ومرج &
\cellcolor{BgFail}E \\
\rowcolor{NeutralLight}
سوریه &
۲۰۱۱– &
چندجانبه &
خیر &
چندگانه &
نیابتی &
۵۰۰K+ &
بی‌نتیجه &
فاجعه &
\cellcolor{BgFail}E–F \\
یمن &
۲۰۱۵– &
عربستان/امارات &
خیر &
دولت هادی &
نظامی &
۳۷۷K+ &
بن‌بست &
بحران انسانی &
\cellcolor{BgFail}E \\
\rowcolor{NeutralLight}
اوکراین &
۲۰۲۲– &
روسیه (مهاجم) &
خیر &
— &
تهاجم &
جاری &
مقاومت اوکراین &
نامشخص &
? \\

\bottomrule
\end{tabular}
\end{adjustbox}
\end{table}
\end{landscape}

% ─────────────────────────────────────────────────────────────────────────
\section{الگوهای استخراج‌شده از مداخلات معاصر}
\label{sec:ch10-patterns}
% ─────────────────────────────────────────────────────────────────────────

\begin{tcolorbox}[wavebox, title={الگوهای کلیدی مداخلات معاصر}]
\begin{enumerate}[nosep, label=\textbf{الگوی \arabic*:}]
  \item \textbf{«سقوط آسان، بازسازی ناممکن»:} سرنگونی نظامی رژیم‌ها آسان بوده (عراق: ۳ هفته، لیبی: ۷ ماه) اما دولت‌سازی پس از آن همواره شکست خورده.

  \item \textbf{مجوز شورای امنیت ≠ موفقیت:} سومالی و لیبی هر دو مجوز داشتند و هر دو شکست خوردند. کوزوو مجوز نداشت و نسبتاً موفق بود. مجوز حقوقی شرط لازم نیست.

  \item \textbf{تأثیر سومالی (Mogadishu Effect):} شکست یک مداخله، اراده‌ی سیاسی برای مداخلات بعدی را تضعیف می‌کند — حتی اگر ضرورت اخلاقی آشکار باشد (رواندا).

  \item \textbf{مداخله‌ی محدود ≠ نتیجه‌ی محدود:} مداخله‌ای که «محدود» آغاز شود (مثلاً «فقط بمباران هوایی» در لیبی) ذاتاً تمایل به گسترش دارد — تأیید الگوی «مست گیرند...».%

  \item \textbf{نجات‌طلبی مؤثر مستلزم ظرفیت داخلی است:} اوکراین (ارتش قوی + دولت مشروع + هویت ملی) در مقایسه با لیبی و سومالی موفق‌تر عمل کرده — زیرا کمک خارجی \concept{مکمل} ظرفیت داخلی بوده نه \concept{جایگزین} آن.

  \item \textbf{گزینشی‌بودن فاحش:} مداخله فقط هنگامی رخ می‌دهد که منافع ژئوپولیتیک ایجاب کند. مقایسه‌ی کوزوو (مداخله) با رواندا (عدم مداخله)، یا سوریه (بن‌بست) با لیبی (مداخله)، این گزینشی‌بودن را آشکار می‌کند.

  \item \textbf{بازگشت منطق جنگ سرد:} پس از ۲۰۱۴ (کریمه)، رقابت ابرقدرتی (امریکا-روسیه-چین) بار دیگر مانع اصلی مداخلات انسان‌دوستانه شده — مانند دوران جنگ سرد.
\end{enumerate}
\end{tcolorbox}

% ─────────────────────────────────────────────────────────────────────────
\section{نمودار تحلیلی: رابطه‌ی ظرفیت داخلی و موفقیت مداخله}
\label{sec:ch10-chart}
% ─────────────────────────────────────────────────────────────────────────

\begin{figure}[htbp]
\centering
\begin{tikzpicture}
\begin{axis}[
  width=0.85\linewidth,
  height=8cm,
  xlabel={\rl{ظرفیت نهادی داخلی (پایین ← بالا)}},
  ylabel={\rl{موفقیت مداخله (F ← A)}},
  xmin=0, xmax=10,
  ymin=0, ymax=10,
  xtick={1,3,5,7,9},
  xticklabels={\rl{بسیار ضعیف},\rl{ضعیف},\rl{متوسط},\rl{قوی},\rl{بسیار قوی}},
  ytick={1,3,5,7,9},
  yticklabels={F,E,D--C,B,A},
  grid=major,
  grid style={dashed, NeutralMid!50},
  axis lines=left,
  every axis label/.style={font=\small},
  every tick label/.style={font=\scriptsize},
]

% نقاط
\node[fail mark] at (axis cs:1,1) {};
\node[right, font=\scriptsize] at (axis cs:1.3,1) {\rl{سومالی}};

\node[fail mark] at (axis cs:1.5,0.5) {};
\node[right, font=\scriptsize] at (axis cs:1.8,0.5) {\rl{رواندا (عدم مداخله)}};

\node[fail mark] at (axis cs:2,2) {};
\node[right, font=\scriptsize] at (axis cs:2.3,2) {\rl{لیبی}};

\node[mixed mark] at (axis cs:3,3.5) {};
\node[right, font=\scriptsize] at (axis cs:3.3,3.5) {\rl{افغانستان}};

\node[mixed mark] at (axis cs:4,4.5) {};
\node[right, font=\scriptsize] at (axis cs:4.3,4.5) {\rl{عراق}};

\node[mixed mark] at (axis cs:2.5,1.5) {};
\node[right, font=\scriptsize] at (axis cs:2.8,1.5) {\rl{سوریه}};

\node[success mark] at (axis cs:5,7) {};
\node[right, font=\scriptsize] at (axis cs:5.3,7) {\rl{کوزوو}};

\node[success mark] at (axis cs:8,9) {};
\node[right, font=\scriptsize] at (axis cs:8.3,9) {\rl{خلیج فارس (کویت)}};

\node[success mark] at (axis cs:9,9.5) {};
\node[right, font=\scriptsize] at (axis cs:7.5,9.8) {\rl{انقلاب شکوهمند}};

\node[mixed mark] at (axis cs:7,7.5) {};
\node[left, font=\scriptsize] at (axis cs:6.7,7.5) {\rl{اوکراین (جاری)}};

\node[success mark] at (axis cs:4.5,7.5) {};
\node[right, font=\scriptsize] at (axis cs:4.8,7.8) {\rl{تیمور شرقی}};

\node[fail mark] at (axis cs:1.5,1.5) {};
\node[left, font=\scriptsize] at (axis cs:1.2,1.8) {\rl{یمن}};

% خط روند
\addplot[domain=0.5:9.5, samples=50, thick, AccentGold, dashed]
  {0.9*x + 0.5};

\node[font=\scriptsize\bfseries, color=AccentGold, rotate=28] at (axis cs:6,6.5)
  {\rl{خط روند: ظرفیت داخلی ↔ موفقیت}};

\end{axis}
\end{tikzpicture}
\caption{رابطه‌ی ظرفیت نهادی داخلی و موفقیت مداخله‌ی خارجی}
\label{fig:ch10-capacity-success}
\end{figure}

% ─────────────────────────────────────────────────────────────────────────
\section{مداخلات منطقه‌ای آفریقایی: اوگاندا، کنگو و ساحل}
\label{sec:ch10-africa}
% ─────────────────────────────────────────────────────────────────────────

\subsection{جنگ‌های کنگو (۱۹۹۶–۲۰۰۳): «جنگ جهانی آفریقا»}

جنگ‌های کنگو — که گاه «جنگ جهانی آفریقا» نامیده می‌شوند — نمونه‌ی بارز \concept{مداخلات منطقه‌ای زنجیره‌ای} هستند: اوگاندا و رواندا از گروه‌های شورشی در کنگو حمایت کردند و سپس خود مستقیماً وارد شدند. نتیجه: بیش از ۵.۴ میلیون کشته — مرگبارترین جنگ پس از جنگ جهانی دوم.%
\footnote{\lr{Stearns, Jason K. \textit{Dancing in the Glory of Monsters: The Collapse of the Congo and the Great War of Africa}. PublicAffairs, 2011.}}

\begin{tcolorbox}[defbox, title={درس آفریقایی}]
جنگ‌های کنگو نشان دادند که:
\begin{itemize}[nosep, rightmargin=1em]
  \item \textbf{مداخله‌ی منطقه‌ای} (نه فقط ابرقدرتی) می‌تواند ویرانگر باشد.
  \item \textbf{«نجات‌طلبی» می‌تواند بهانه‌ی غارت باشد:} اوگاندا و رواندا از منابع معدنی کنگو (کلتان، الماس) بهره بردند.%
  \footnote{\lr{UN Panel of Experts. "Illegal Exploitation of Natural Resources in the DRC." S/2001/357, 2001.}}
  \item \textbf{فقدان توجه بین‌المللی:} با وجود مرگبارترین جنگ معاصر، کنگو هرگز توجه رسانه‌ای و سیاسی مشابه کوزوو یا عراق را دریافت نکرد — تأیید گزینشی‌بودن حمایت بین‌المللی.
\end{itemize}
\end{tcolorbox}

\subsection{اوگاندا: کودتاهای پیاپی و مداخلات}

تاریخ اوگاندا مملو از مداخلات خارجی و نجات‌طلبی است:

\begin{longtable}{
  >{\raggedleft\arraybackslash\bfseries}p{2.5cm}
  >{\raggedleft\arraybackslash}p{2cm}
  >{\raggedleft\arraybackslash}p{3cm}
  >{\raggedleft\arraybackslash}p{2.5cm}
  >{\raggedleft\arraybackslash}p{3cm}}
\caption{مداخلات و نجات‌طلبی در اوگاندا}\label{tab:ch10-uganda}\\
\toprule
\rowcolor{PrimaryDeep}
\textcolor{white}{\textbf{رویداد}} &
\textcolor{white}{\textbf{سال}} &
\textcolor{white}{\textbf{مداخله‌گر/حامی}} &
\textcolor{white}{\textbf{ماهیت}} &
\textcolor{white}{\textbf{نتیجه}} \\
\midrule
\endfirsthead
\bottomrule
\endlastfoot

سقوط ایدی امین &
۱۹۷۹ &
تانزانیا + تبعیدیان اوگاندایی &
نظامی مستقیم &
سقوط امین؛ بی‌ثباتی ادامه%
\footnote{\lr{Avirgan, Tony, and Martha Honey. \textit{War in Uganda: The Legacy of Idi Amin}. Lawrence Hill, 1982.}} \\
\rowcolor{NeutralLight}
جنگ بوش (موسِوِنی) &
۱۹۸۱–۱۹۸۶ &
حمایت لیبی و کنیا (محدود) &
شورش مسلحانه &
پیروزی موسِوِنی \\
مداخله در کنگو &
۱۹۹۶–۲۰۰۳ &
اوگاندا (به‌عنوان مداخله‌گر) &
«نجات» + غارت &
بی‌ثباتی منطقه‌ای \\

\end{longtable}

% ─────────────────────────────────────────────────────────────────────────
\section{یمن (۲۰۱۵–کنون): مداخله‌ی عربستان و بحران انسانی}
\label{sec:ch10-yemen}
% ─────────────────────────────────────────────────────────────────────────

\subsection{زمینه و تحلیل}

مداخله‌ی نظامی ائتلاف به رهبری عربستان سعودی در یمن (مارس ۲۰۱۵) به «دعوت» \histref{عبدربه منصور هادی}{رئیس‌جمهور مستعفی یمن} صورت گرفت. هادی از عربستان خواست تا علیه \enemy{حوثی‌ها} (انصارالله) مداخله نظامی کند.%
\footnote{\lr{Salisbury, Peter. "Yemen: National Chaos, Local Order." Chatham House Report, December 2017.}}

\begin{tcolorbox}[casebox, title={تحلیل شش‌وجهی: یمن (۲۰۱۵)}]
\begin{itemize}[nosep, rightmargin=1em]
  \item \textbf{ساختاری:} یمن پیش از مداخله دولت شکننده‌ای بود. هادی مشروعیت محدودی داشت (دوره‌ی ریاست‌جمهوری‌اش منقضی شده بود).
  \item \textbf{حقوقی:} بدون مجوز شورای امنیت (قطعنامه ۲۲۱۶ تحریم تسلیحاتی حوثی‌ها را تصویب کرد اما مجوز مداخله‌ی نظامی مستقیم نداد). ادعای «دعوت حکومت مشروع» از نظر حقوقی بحث‌برانگیز.%
  \footnote{\lr{Ruys and Ferro, 2016.}}
  \item \textbf{اقتصادی:} یمن فقیرترین کشور عربی. عربستان نگران نفوذ ایران بود نه منافع اقتصادی.
  \item \textbf{نظامی:} برتری هوایی مطلق ائتلاف اما ناتوانی در پیشروی زمینی. حوثی‌ها از حمایت ایرانی (تسلیحات و مشاوره) بهره بردند.
  \item \textbf{روان‌شناختی:} رقابت فرقه‌ای (سنّی-شیعه) و ترس از «هلال شیعی» ایران.
  \item \textbf{ژئوپولیتیک:} جنگ نیابتی عربستان-ایران. امریکا و بریتانیا با فروش سلاح به ائتلاف عملاً شریک بودند.%
  \footnote{\lr{Riedel, Bruce. "Who Are the Houthis, and Why Are We at War with Them?" Brookings, December 2017.}}
  \item \textbf{نتیجه (طیف):} \textbf{E}. بیش از ۱۵۰٬۰۰۰ کشته‌ی مستقیم و بیش از ۲۲۷٬۰۰۰ مرگ ناشی از قحطی و بیماری. «بدترین بحران انسانی جهان» (سازمان ملل).%
  \footnote{\lr{UN OCHA. "Yemen Humanitarian Crisis Overview." 2023.}}
\end{itemize}
\end{tcolorbox}

% ─────────────────────────────────────────────────────────────────────────
\section{نمودار مقایسه‌ای: هزینه‌های انسانی مداخلات}
\label{sec:ch10-human-cost}
% ─────────────────────────────────────────────────────────────────────────

\begin{figure}[htbp]
\centering
\begin{tikzpicture}
\begin{axis}[
  xbar,
  width=0.85\linewidth,
  height=10cm,
  xlabel={\rl{تعداد کشته‌ها (هزار نفر)}},
  symbolic y coords={
    {\rl{خلیج فارس}},
    {\rl{کوزوو}},
    {\rl{سومالی}},
    {\rl{تیمور شرقی}},
    {\rl{افغانستان}},
    {\rl{یمن}},
    {\rl{عراق}},
    {\rl{سوریه}},
    {\rl{رواندا}},
    {\rl{کنگو}}
  },
  ytick=data,
  nodes near coords,
  nodes near coords align={horizontal},
  every node near coord/.style={font=\scriptsize},
  bar width=12pt,
  xmin=0, xmax=6000,
  axis lines=left,
  every axis label/.style={font=\small},
  every tick label/.style={font=\scriptsize},
  enlarge y limits=0.08,
]

\addplot[fill=AccentGreen!60, draw=AccentGreen] coordinates {
  (1,{\rl{خلیج فارس}})
  (13,{\rl{کوزوو}})
};

\addplot[fill=AccentAmber!60, draw=AccentAmber] coordinates {
  (10,{\rl{تیمور شرقی}})
  (30,{\rl{سومالی}})
  (170,{\rl{افغانستان}})
};

\addplot[fill=AccentRed!60, draw=AccentRed] coordinates {
  (200,{\rl{عراق}})
  (377,{\rl{یمن}})
  (500,{\rl{سوریه}})
  (800,{\rl{رواندا}})
  (5400,{\rl{کنگو}})
};

\end{axis}
\end{tikzpicture}
\caption{مقایسه‌ی هزینه‌های انسانی مداخلات معاصر (هزار کشته)}
\label{fig:ch10-human-cost}
\end{figure}

% ─────────────────────────────────────────────────────────────────────────
\section{تحلیل ترکیبی: چرا برخی مداخلات «موفق‌تر» بوده‌اند؟}
\label{sec:ch10-success-factors}
% ─────────────────────────────────────────────────────────────────────────

بر مبنای بررسی تطبیقی تمامی موارد این فصل:

\begin{tcolorbox}[successbox, title={عوامل مشترک مداخلات نسبتاً موفق (A–B)}]
\begin{enumerate}[nosep, label=\textbf{\arabic*.}]
  \item \textbf{هدف محدود و مشخص:} آزادی کویت (نه تغییر رژیم عراق)؛ توقف پاک‌سازی قومی کوزوو (نه تغییر رژیم صربستان).
  \item \textbf{ظرفیت نهادی داخلی:} کویت دولت مشخصی داشت؛ کوزوو نهادهای در حال شکل‌گیری داشت؛ اوکراین ارتش و دولت قوی دارد.
  \item \textbf{حمایت گسترده‌ی بین‌المللی:} خلیج فارس (۳۴ کشور)؛ اوکراین (حمایت وسیع غربی).
  \item \textbf{استراتژی خروج:} خلیج فارس (خروج پس از آزادی)؛ انقلاب شکوهمند (ویلیام پادشاه مشروطه شد).
  \item \textbf{عدم تلاش برای «دولت‌سازی» از صفر:} مداخلاتی که سعی در بازآفرینی کامل جامعه نکردند، موفق‌تر بودند.
\end{enumerate}
\end{tcolorbox}

\begin{tcolorbox}[failbox, title={عوامل مشترک مداخلات شکست‌خورده (D–F)}]
\begin{enumerate}[nosep, label=\textbf{\arabic*.}]
  \item \textbf{هدف مبهم و متغیر:} افغانستان (از «نابودی القاعده» به «دموکراسی‌سازی»)؛ لیبی (از «حمایت» به «تغییر رژیم»).
  \item \textbf{فقدان ظرفیت نهادی:} سومالی، لیبی، یمن — بدون نهاد داخلی قابل‌اتکا.
  \item \textbf{آرزواندیشی مداخله‌گر:} باور به امکان تبدیل جوامع قبیله‌ای/فرقه‌ای به دموکراسی لیبرال در مدت کوتاه.%
  \footnote{\lr{SIGAR, 2021. "What We Need to Learn: Lessons from Twenty Years of Afghanistan Reconstruction."}}
  \item \textbf{فقدان استراتژی خروج:} افغانستان (۲۰ سال)؛ عراق (۸ سال + بازگشت).
  \item \textbf{مخالفت منطقه‌ای:} پاکستان (افغانستان)؛ ایران (عراق)؛ روسیه (سوریه).
  \item \textbf{تغییر هدف در حین عملیات} (\lr{mission creep}): لیبی مهم‌ترین نمونه.
\end{enumerate}
\end{tcolorbox}

% ─────────────────────────────────────────────────────────────────────────
\section{چشم‌انداز آینده: نجات‌طلبی در عصر رقابت ابرقدرتی}
\label{sec:ch10-future}
% ─────────────────────────────────────────────────────────────────────────

\begin{tcolorbox}[defbox, title={روندهای آینده}]
\begin{enumerate}[nosep, label=\textbf{\arabic*.}]
  \item \textbf{بازگشت وتو:} رقابت امریکا-روسیه-چین مانع هرگونه مجوز شورای امنیت برای مداخله‌ی جدید خواهد شد — مگر در بحران‌هایی که هیچ ابرقدرتی ذی‌نفع نباشد.

  \item \textbf{مداخلات منطقه‌ای:} قدرت‌های منطقه‌ای (ترکیه، عربستان، امارات، ایران، اتیوپی) به‌جای ابرقدرت‌ها مداخله خواهند کرد — اما با منابع کمتر و پاسخ‌گویی کمتر.

  \item \textbf{جنگ سایبری و اطلاعاتی:} «مداخله» دیگر الزاماً نظامی نیست. دخالت در انتخابات، جنگ سایبری و تبلیغات رسانه‌ای، ابزارهای نوین «نجات‌طلبی» و «مداخله» هستند.%
  \footnote{\lr{Rid, Thomas. \textit{Active Measures: The Secret History of Disinformation and Political Warfare}. Farrar, Straus and Giroux, 2020.}}

  \item \textbf{تغییرات اقلیمی:} بحران‌های ناشی از تغییر اقلیم (خشکسالی، سیل، مهاجرت اجباری) نوع جدیدی از «بحران انسان‌دوستانه» ایجاد خواهد کرد که مداخله‌ی بین‌المللی را ضروری اما دشوارتر می‌سازد.

  \item \textbf{فناوری و نجات‌طلبی:} شبکه‌های اجتماعی امکان «نجات‌طلبی» مردمی را تسهیل کرده (بهار عربی) اما هم‌زمان ابزار «پروپاگاندای مداخله‌گر» نیز شده‌اند.
\end{enumerate}
\end{tcolorbox}

% ─────────────────────────────────────────────────────────────────────────
\section{جمع‌بندی فصل}
\label{sec:ch10-summary}
% ─────────────────────────────────────────────────────────────────────────

\begin{tcolorbox}[wavebox, title={یافته‌های کلیدی فصل ۱۰}]
\begin{enumerate}[nosep, label=\textbf{\arabic*.}]
  \item در ۵۰ سال اخیر، \textbf{هیچ مداخله‌ی نظامی‌ای که هدفش «دولت‌سازی از صفر» بوده، موفق نشده} (افغانستان، عراق، لیبی، سومالی).

  \item مداخلاتی نسبتاً موفق بوده‌اند که \textbf{هدف محدود و مشخص، ظرفیت داخلی قابل‌اتکا، و استراتژی خروج} داشته‌اند (کویت، کوزوو، تیمور شرقی).

  \item \textbf{عدم مداخله نیز فاجعه‌بار است:} رواندا نشان داد که «بی‌عملی» می‌تواند ۸۰۰٬۰۰۰ جان ببرد.

  \item \textbf{مجوز شورای امنیت لازم اما کافی نیست:} سومالی و لیبی مجوز داشتند و شکست خوردند.

  \item \textbf{گزینشی‌بودن} بزرگ‌ترین ضعف اخلاقی نظام بین‌الملل است: مقایسه‌ی واکنش به کوزوو (مداخله) با رواندا (بی‌عملی) و کنگو (بی‌تفاوتی) این ضعف را آشکار می‌کند.

  \item \textbf{رابطه‌ی خطی} میان ظرفیت نهادی داخلی و موفقیت مداخله وجود دارد (شکل~\ref{fig:ch10-capacity-success}).

  \item \textbf{آینده‌ی تاریک:} رقابت ابرقدرتی، تغییرات اقلیمی و جنگ سایبری، مداخلات انسان‌دوستانه را در آینده دشوارتر خواهد کرد.

  \item \textbf{قاعده‌ی طلایی همچنان معتبر است:} «مست گیرند، در شهر هر آنکه هست گیرند» — هر مداخله‌ای، حتی با بهترین نیت‌ها، ذاتاً گسترش‌یابنده و غیرقابل‌کنترل است. تنها مداخلاتی نسبتاً موفق بوده‌اند که \textbf{مکمل} ظرفیت داخلی بوده‌اند نه \textbf{جایگزین} آن.
\end{enumerate}
\end{tcolorbox}

\cleardoublepage